% GENERATED FILE. Edit canonical lessons, then run npm run generate:books.
% canonical-source-hash: fnv1a64:dd47fec48de19a45
% canonical-lessons: SA-C500-devarah, SA-C500-nananda, SA-C500-syalah, SA-C500-prapitamahah, SA-C500-santatih

\chapter{In-Laws and Ancestors}
\label{ch:sa-in-laws-and-ancestors}
\hlchaptermodality{500}

\begin{chapteropening}
\textbf{By the end of this chapter:} \emph{I can name my in-laws and my ancestors.}

Complete the last lesson of chapter 500: I can name my in-laws and my ancestors.
\end{chapteropening}

\section[\texorpdfstring{\sk{देवरः} (devaraḥ)}{देवरः (devaraḥ)}]{\sk{देवरः} (devaraḥ) — a husband's brother}

\label{lesson:SA-C500-devarah}



\begin{quote}
Before the new one: say the Sanskrit for a sister's son, then the Sanskrit for a daughter's son.
\end{quote}

\subsection*{You'll want to know: \sk{देवरः}}
\textbf{\sk{देवरः}} — \emph{devaraḥ} — \textquotedblleft{}a husband's brother\textquotedblright{}.

\begin{quote}
\sk{मम} \sk{देवरः} \sk{चिकित्सकः} \sk{अस्ति।}
\end{quote}

\begin{quote}
\emph{mama devaraḥ cikitsakaḥ asti.}
\end{quote}

\begin{quote}
My husband's brother is a doctor.
\end{quote}

A masculine noun.

\begin{cousinweb}
\textbf{\sk{देवरः}} (\emph{devaraḥ}) is built on \emph{devṛ}, \textquotedblleft{}a husband's brother (esp. his younger brother)\textquotedblright{}; MW compares Greek \emph{daēr} and Latin \emph{levir} (MW).
\end{cousinweb}

\begin{grammarlens}[title={What you've built}]
One more word for describing what you lived and felt.
\end{grammarlens}

\subsection*{Your turn}
\begin{itemize}
\raggedright
  \item \emph{Say it:} \emph{devaraḥ}
  \item \emph{Say it:} \emph{devaraḥ}, once more
  \item \emph{Say it:} \emph{mama devaraḥ cikitsakaḥ asti.} — My husband's brother is a doctor.
  \item \emph{Recall:} say the Sanskrit for a sister's son, then the Sanskrit for a daughter's son, then say \emph{devaraḥ} again
\end{itemize}

\begin{tcolorbox}[breakable,colback=teal!4,colframe=teal!35!black,title={Before you move on}]
What does \sk{देवरः} mean? (\textquotedblleft{}A husband's brother\textquotedblright{}.) Say it once more.
\end{tcolorbox}
\section[\texorpdfstring{\sk{ननान्दा} (nanāndā)}{ननान्दा (nanāndā)}]{\sk{ननान्दा} (nanāndā) — a husband's sister}

\label{lesson:SA-C500-nananda}



\begin{quote}
Before the new one: say the Sanskrit for a daughter's son, then the Sanskrit for a husband's brother.
\end{quote}

\subsection*{You'll want to know: \sk{ननान्दा}}
\textbf{\sk{ननान्दा}} — \emph{nanāndā} — \textquotedblleft{}a husband's sister\textquotedblright{}.

\begin{quote}
\sk{ननान्दा} \sk{मया} \sk{सह} \sk{आपणं} \sk{गच्छति।}
\end{quote}

\begin{quote}
\emph{nanāndā mayā saha āpaṇaṃ gacchati.}
\end{quote}

\begin{quote}
My husband's sister goes to the shop with me.
\end{quote}

A feminine noun.

\begin{cousinweb}
\textbf{\sk{ननान्दा}} (\emph{nanāndā}) MW records \emph{nanāndṛ} (Rigveda) and \emph{nanandṛ}, \textquotedblleft{}a husband's sister\textquotedblright{}; MW gives no derivation, so its origin is uncertain (MW).
\end{cousinweb}

\begin{grammarlens}[title={What you've built}]
One more word for describing what you lived and felt.
\end{grammarlens}

\subsection*{Your turn}
\begin{itemize}
\raggedright
  \item \emph{Say it:} \emph{nanāndā}
  \item \emph{Say it:} \emph{nanāndā}, once more
  \item \emph{Say it:} \emph{nanāndā mayā saha āpaṇaṃ gacchati.} — My husband's sister goes to the shop with me.
  \item \emph{Recall:} say the Sanskrit for a daughter's son, then the Sanskrit for a husband's brother, then say \emph{nanāndā} again
\end{itemize}

\begin{tcolorbox}[breakable,colback=teal!4,colframe=teal!35!black,title={Before you move on}]
What does \sk{ननान्दा} mean? (\textquotedblleft{}A husband's sister\textquotedblright{}.) Say it once more.
\end{tcolorbox}
\section[\texorpdfstring{\sk{स्यालः} (syālaḥ)}{स्यालः (syālaḥ)}]{\sk{स्यालः} (syālaḥ) — a wife's brother}

\label{lesson:SA-C500-syalah}



\begin{quote}
Before the new one: say the Sanskrit for a husband's brother, then the Sanskrit for a husband's sister.
\end{quote}

\subsection*{You'll want to know: \sk{स्यालः}}
\textbf{\sk{स्यालः}} — \emph{syālaḥ} — \textquotedblleft{}a wife's brother\textquotedblright{}.

\begin{quote}
\sk{मम} \sk{स्यालः} \sk{अद्य} \sk{आगच्छति।}
\end{quote}

\begin{quote}
\emph{mama syālaḥ adya āgacchati.}
\end{quote}

\begin{quote}
My wife's brother is coming today.
\end{quote}

A masculine noun.

\begin{cousinweb}
\textbf{\sk{स्यालः}} (\emph{syālaḥ}) MW calls it of doubtful derivation, so its origin is uncertain; it gives \textquotedblleft{}a wife's brother\textquotedblright{} and notes the spelling \emph{śyāla} (MW).
\end{cousinweb}

\begin{grammarlens}[title={What you've built}]
One more word for describing what you lived and felt.
\end{grammarlens}

\subsection*{Your turn}
\begin{itemize}
\raggedright
  \item \emph{Say it:} \emph{syālaḥ}
  \item \emph{Say it:} \emph{syālaḥ}, once more
  \item \emph{Say it:} \emph{mama syālaḥ adya āgacchati.} — My wife's brother is coming today.
  \item \emph{Recall:} say the Sanskrit for a husband's brother, then the Sanskrit for a husband's sister, then say \emph{syālaḥ} again
\end{itemize}

\begin{tcolorbox}[breakable,colback=teal!4,colframe=teal!35!black,title={Before you move on}]
What does \sk{स्यालः} mean? (\textquotedblleft{}A wife's brother\textquotedblright{}.) Say it once more.
\end{tcolorbox}
\section[\texorpdfstring{\sk{प्रपितामहः} (prapitāmahaḥ)}{प्रपितामहः (prapitāmahaḥ)}]{\sk{प्रपितामहः} (prapitāmahaḥ) — a great-grandfather}

\label{lesson:SA-C500-prapitamahah}



\begin{quote}
Before the new one: say the Sanskrit for a husband's sister, then the Sanskrit for a wife's brother.
\end{quote}

\subsection*{You'll want to know: \sk{प्रपितामहः}}
\textbf{\sk{प्रपितामहः}} — \emph{prapitāmahaḥ} — \textquotedblleft{}a great-grandfather\textquotedblright{}.

\begin{quote}
\sk{मम} \sk{प्रपितामहः} \sk{कृषकः} \sk{अभवत्।}
\end{quote}

\begin{quote}
\emph{mama prapitāmahaḥ kṛṣakaḥ abhavat.}
\end{quote}

\begin{quote}
My great-grandfather was a farmer.
\end{quote}

A masculine noun.

\begin{cousinweb}
\textbf{\sk{प्रपितामहः}} (\emph{prapitāmahaḥ}) is built from \emph{pra}, \textquotedblleft{}before\textquotedblright{}, and \textbf{\sk{पितामहः}}, \textquotedblleft{}a grandfather\textquotedblright{}; MW gives \textquotedblleft{}a paternal great-grandfather\textquotedblright{} and, in the plural, \textquotedblleft{}ancestors\textquotedblright{} (MW).
\end{cousinweb}

\begin{grammarlens}[title={What you've built}]
One more word for describing what you lived and felt.
\end{grammarlens}

\subsection*{Your turn}
\begin{itemize}
\raggedright
  \item \emph{Say it:} \emph{prapitāmahaḥ}
  \item \emph{Say it:} \emph{prapitāmahaḥ}, once more
  \item \emph{Say it:} \emph{mama prapitāmahaḥ kṛṣakaḥ abhavat.} — My great-grandfather was a farmer.
  \item \emph{Recall:} say the Sanskrit for a husband's sister, then the Sanskrit for a wife's brother, then say \emph{prapitāmahaḥ} again
\end{itemize}

\begin{tcolorbox}[breakable,colback=teal!4,colframe=teal!35!black,title={Before you move on}]
What does \sk{प्रपितामहः} mean? (\textquotedblleft{}A great-grandfather\textquotedblright{}.) Say it once more.
\end{tcolorbox}
\section[\texorpdfstring{\sk{सन्ततिः} (santatiḥ)}{सन्ततिः (santatiḥ)}]{\sk{सन्ततिः} (santatiḥ) — offspring}

\label{lesson:SA-C500-santatih}



\begin{quote}
Before the new one: say the Sanskrit for a wife's brother, then the Sanskrit for a great-grandfather.
\end{quote}

\subsection*{You'll want to know: \sk{सन्ततिः}}
\textbf{\sk{सन्ततिः}} — \emph{santatiḥ} — \textquotedblleft{}offspring\textquotedblright{}.

\begin{quote}
\sk{माता} \sk{पिता} \sk{च} \sk{सन्ततेः} \sk{सुखं} \sk{चिन्तयतः।}
\end{quote}

\begin{quote}
\emph{mātā pitā ca santateḥ sukhaṃ cintayataḥ.}
\end{quote}

\begin{quote}
Parents think of their children's happiness.
\end{quote}

A feminine noun.

\begin{cousinweb}
\textbf{\sk{सन्ततिः}} (\emph{santatiḥ}) is built from \emph{sam} and the root √tan, \textquotedblleft{}to stretch\textquotedblright{}: a line stretched on; MW gives \textquotedblleft{}uninterrupted succession, lineage, progeny, offspring\textquotedblright{} (MW).
\end{cousinweb}

\begin{grammarlens}[title={What you've built}]
One more word for describing what you lived and felt.
\end{grammarlens}

\subsection*{Your turn}
\begin{itemize}
\raggedright
  \item \emph{Say it:} \emph{santatiḥ}
  \item \emph{Say it:} \emph{santatiḥ}, once more
  \item \emph{Say it:} \emph{mātā pitā ca santateḥ sukhaṃ cintayataḥ.} — Parents think of their children's happiness.
  \item \emph{Recall:} say the Sanskrit for a wife's brother, then the Sanskrit for a great-grandfather, then say \emph{santatiḥ} again
\end{itemize}

\begin{tcolorbox}[breakable,colback=teal!4,colframe=teal!35!black,title={Before you move on}]
What does \sk{सन्ततिः} mean? (\textquotedblleft{}Offspring\textquotedblright{}.) Say it once more.
\end{tcolorbox}
